% defined-in: apostol (page 174)
% entity-id: thm-weierstrass-extreme-value
% entity-type: theorem

\entityref{op-hypertarget}{\hypertarget}{thm-weierstrass-extreme-value}{}
\hypertarget{thm-weierstrass-extreme-value}{}
\begin{theorem}[Weierstrass Extreme Value]
\[
\mForall{\entityref{obj-function}{f} \colon [a, b] \mTo \entityref{obj-real-numbers}{\mathbb{R}}}
\quad \entityref{prop-continuous}{f} \mImplies \mExists{c, d \in [a, b]} \quad f(c) = \entityref{op-supremum}{\sup\limits_{x \in [a, b]}} f(x) \entityref{term-conjunction}{\land} f(d) = \entityref{op-infimum}{\inf\limits_{x \in [a, b]}} f(x)
\]
\end{theorem}

\begin{proof}[RU]
Достаточно доказать, что \entityref{prop-continuous}{f} достигает своей верхней грани на $[a, b]$. Пусть $M = \entityref{op-supremum}{\sup\limits_{x \in [a, b]}} f(x)$. Предположим, что не существует такого $x \in [a, b]$, для которого $f(x) = M$. Определим $g(x) = M - f(x)$. Тогда $g(x) > 0$ для всех $x \in [a, b]$, следовательно, \entityref{obj-function}{функция} $1/\entityref{obj-function}{g}$ непрерывна на $[a, b]$. Согласно теореме о ограниченности непрерывной \entityref{obj-function}{функции}, $1/\entityref{obj-function}{g}$ ограничена на $[a, b]$, то есть $1/g(x) < C$ для некоторого $C > 0$. Это влечет $M - f(x) > 1/C$, откуда $f(x) < M - 1/C$ для всех $x \in [a, b]$. Это противоречит тому, что $M$ является точной верхней гранью $\entityref{obj-function}{f}$ на $[a, b]$. Следовательно, $f(x) = M$ для некоторого $x \in [a, b]$. Аналогичное рассуждение применимо к $\entityref{op-infimum}{\inf} \entityref{obj-function}{f}$.
\end{proof}

\begin{proof}[EN]
It suffices to prove that \entityref{prop-continuous}{f} attains its supremum on $[a, b]$. Let $M = \entityref{op-supremum}{\sup\limits_{x \in [a, b]}} f(x)$. Assume there is no $x \in [a, b]$ such that $f(x) = M$. Define $g(x) = M - f(x)$. Then $g(x) > 0$ for all $x \in [a, b]$, so the reciprocal $1/\entityref{obj-function}{g}$ is continuous on $[a, b]$. By the theorem on the boundedness of continuous functions, $1/\entityref{obj-function}{g}$ is bounded on $[a, b]$, such that $1/g(x) < C$ for some $C > 0$. This implies $M - f(x) > 1/C$, so that $f(x) < M - 1/C$ for all $x \in [a, b]$. This contradicts the fact that $M$ is the least upper bound of $\entityref{obj-function}{f}$ on $[a, b]$. Hence, $f(x) = M$ for at least one $x \in [a, b]$. A similar argument applies to $\entityref{op-infimum}{\inf} \entityref{obj-function}{f}$.
\end{proof}